\chapter{Observability, operations, incidents, and postmortems}

\section{From running to knowing}
A running process is not necessarily a useful service. Operations asks whether
the service is delivering its intended outcome and what a responsible person can
do when it is not. Observability is the ability to infer internal state from
externally emitted signals. OpenTelemetry defines a vendor-neutral set of
telemetry concepts and uses normative terms such as MUST and SHOULD inside its
specification \cite{opentelemetry2026spec}.

\section{Signals}
\term{Metrics} summarize values over time: request count, error count, latency,
queue age, saturation. \term{Logs} record discrete events with context.
\term{Traces} connect work across process and network boundaries. Use each for
the question it answers. A trace is not a replacement for a durable audit log,
and a metric cannot show every cause.

Every signal needs a schema and a cost model. Include request or trace identity,
route, outcome, dependency, and version where useful. Avoid high-cardinality
labels such as raw user IDs in metrics. Redact secrets and personal data before
emission, not only in a dashboard.

\section{Alerting}
An alert should indicate an actionable condition, have an owner, and explain the
first safe action. Alert on symptoms that matter to users or operators: an SLO
burn rate, queue age beyond a business limit, or inability to accept an order.
Avoid alerts that merely say a process is busy unless somebody can act on them.

A dashboard is a decision surface. Arrange it from outcome to cause: service
health, user journey, traffic, errors, latency, saturation, dependencies, and
recent changes. Test it during a failure; a beautiful dashboard that cannot
answer “what changed?” is not operationally useful.

\section{Runbooks and safe operations}
A runbook should say when it applies, what evidence to collect, what actions are
safe, what actions are irreversible, how to escalate, and how to verify recovery.
Prefer reversible controls: reduce traffic, disable a feature, pause a consumer,
or increase capacity. Administrative endpoints need authentication,
authorization, audit, and a break-glass procedure.

\section{Incident response}
During an incident, stabilize first. Establish a coordinator, a communications
channel, an incident timeline, and an explicit hypothesis. Separate observations
from guesses. Record commands and configuration changes. Avoid risky experiments
on a system whose state is not understood.

An incident is a period where the system's expected behaviour is not being met;
it is not limited to outages. A privacy exposure, incorrect financial result, or
unreliable deployment can be an incident even when servers respond normally.

\section{Postmortems}
A useful postmortem is blameless about individual intent and precise about system
conditions. Include impact, detection, timeline, contributing factors, what
worked, what failed, and actions with owners and dates. Do not produce only a
long list of “be more careful.” Improve a control: add a validation, limit blast
radius, make a signal actionable, rehearse recovery, or change an approval path.

\section{Change and observability}
Observability must evolve with the system. A new queue requires queue age and
dead-letter signals. A new authorization boundary requires denied-access audit
events. A new batch job needs freshness and completeness, not only process
uptime. Treat telemetry schemas as compatibility-sensitive APIs.

\begin{exercise}
Design an alert for an order service whose success SLO is burning budget rapidly.
Name the indicator, aggregation window, owner, first action, and one false
positive. Then write the minimum runbook needed to pause new intake safely.
\end{exercise}

\section{Chapter summary}
Operations is part of the product. Emit useful, privacy-aware signals; alert on
actionable user impact; maintain reversible runbooks; coordinate incidents with
evidence; and convert postmortem learning into stronger system controls.

